% GENERATED FILE. Edit canonical lessons, then run npm run generate:books.
% canonical-source-hash: fnv1a64:b7759939a53a9b05
% canonical-lessons: HI-C516-tapkana, HI-C516-darakna, HI-C516-paip, HI-C516-ched, HI-C516-dhalakna

\chapter{A Leak in the Roof}
\label{ch:hi-a-leak-in-the-roof}
\hlchaptermodality{516}

\begin{chapteropening}
\textbf{By the end of this chapter:} \emph{I can describe leaks, damp and cracks at home.}

Complete the last lesson of chapter 516: I can describe leaks, damp and cracks at home.
\end{chapteropening}

\section[\texorpdfstring{\hi{टपकाना} (ṭapkānā)}{टपकाना (ṭapkānā)}]{\hi{टपकाना} (ṭapkānā) — to let drip; to drip (something) in}

\label{lesson:HI-C516-tapkana}



\begin{quote}
Before the new one: say the Hindi for to tighten, then the Hindi for to rattle.
\end{quote}

\subsection*{You'll want to know: \hi{टपकाना}}
\textbf{\hi{टपकाना}} — \emph{ṭapkānā} — \textquotedblleft{}to let drip; to drip (something) in\textquotedblright{}.

\begin{quote}
\hi{पुराना} \hi{नल} \hi{रात} \hi{भर} \hi{पानी} \hi{टपकाता} \hi{है।}
\end{quote}

\begin{quote}
\emph{Purānā nal rāt bhar pānī ṭapkātā hai.}
\end{quote}

\begin{quote}
The old tap drips water all night.
\end{quote}

Past, with \hi{ने} (he or she …): \hi{उसने} \hi{आँख} \hi{में} \hi{दवा} \hi{टपकाई}, \emph{usne āṁkh meṁ davā ṭapkāī}.

\begin{cousinweb}
\textbf{\hi{टपकाना}} (\emph{ṭapkānā}) is the causative of \emph{ṭapaknā}, \textquotedblleft{}to drip\textquotedblright{}, which CDIAL files under \emph{*ṭapp}, a reconstructed form (CDIAL 5444).
\end{cousinweb}

\begin{grammarlens}[title={What you've built}]
One more word for describing what you lived and felt.
\end{grammarlens}

\subsection*{Your turn}
\begin{itemize}
\raggedright
  \item \emph{Say it:} \emph{ṭapkānā}
  \item \emph{Say it:} \emph{ṭapkānā}, once more
  \item \emph{Say it:} \emph{Purānā nal rāt bhar pānī ṭapkātā hai.} — The old tap drips water all night.
  \item \emph{Recall:} say the Hindi for to tighten, then the Hindi for to rattle, then say \emph{ṭapkānā} again
\end{itemize}

\begin{tcolorbox}[breakable,colback=teal!4,colframe=teal!35!black,title={Before you move on}]
What does \hi{टपकाना} mean? (\textquotedblleft{}To let drip; to drip (something) in\textquotedblright{}.) Say it once more.
\end{tcolorbox}
\section[\texorpdfstring{\hi{दरकना} (daraknā)}{दरकना (daraknā)}]{\hi{दरकना} (daraknā) — to crack, to split}

\label{lesson:HI-C516-darakna}



\begin{quote}
Before the new one: say the Hindi for to rattle, then the Hindi for to let drip.
\end{quote}

\subsection*{You'll want to know: \hi{दरकना}}
\textbf{\hi{दरकना}} — \emph{daraknā} — \textquotedblleft{}to crack, to split\textquotedblright{}.

\begin{quote}
\hi{भूकंप} \hi{के} \hi{बाद} \hi{दीवार} \hi{दरक} \hi{गई।}
\end{quote}

\begin{quote}
\emph{Bhūkamp ke bād dīvār darak gaī.}
\end{quote}

\begin{quote}
After the earthquake the wall cracked.
\end{quote}

Past (it …): \hi{वह} \hi{शीशा} \hi{दरक} \hi{गया}, \emph{vah śīśā darak gayā}.

\begin{cousinweb}
\textbf{\hi{दरकना}} (\emph{daraknā}) is a tadbhava, inherited from Sanskrit \emph{darayati}, \textquotedblleft{}causes to split\textquotedblright{} (CDIAL 6192). CDIAL gives the Hindi verb as \textquotedblleft{}to be torn\textquotedblright{}.
\end{cousinweb}

\begin{grammarlens}[title={What you've built}]
One more word for describing what you lived and felt.
\end{grammarlens}

\subsection*{Your turn}
\begin{itemize}
\raggedright
  \item \emph{Say it:} \emph{daraknā}
  \item \emph{Say it:} \emph{daraknā}, once more
  \item \emph{Say it:} \emph{Bhūkamp ke bād dīvār darak gaī.} — After the earthquake the wall cracked.
  \item \emph{Recall:} say the Hindi for to rattle, then the Hindi for to let drip, then say \emph{daraknā} again
\end{itemize}

\begin{tcolorbox}[breakable,colback=teal!4,colframe=teal!35!black,title={Before you move on}]
What does \hi{दरकना} mean? (\textquotedblleft{}To crack, to split\textquotedblright{}.) Say it once more.
\end{tcolorbox}
\section[\texorpdfstring{\hi{पाइप} (pāip)}{पाइप (pāip)}]{\hi{पाइप} (pāip) — a pipe}

\label{lesson:HI-C516-paip}



\begin{quote}
Before the new one: say the Hindi for to let drip, then the Hindi for to crack, to split.
\end{quote}

\subsection*{You'll want to know: \hi{पाइप}}
\textbf{\hi{पाइप}} — \emph{pāip} — \textquotedblleft{}a pipe\textquotedblright{}.

\begin{quote}
\hi{रसोई} \hi{का} \hi{पाइप} \hi{टूट} \hi{गया।}
\end{quote}

\begin{quote}
\emph{Rasoī kā pāip ṭūṭ gayā.}
\end{quote}

\begin{quote}
The kitchen pipe has broken.
\end{quote}

Masculine.

\begin{cousinweb}
\textbf{\hi{पाइप}} (\emph{pāip}) is a loan from English \emph{pipe}.
\end{cousinweb}

\begin{grammarlens}[title={What you've built}]
One more word for describing what you lived and felt.
\end{grammarlens}

\subsection*{Your turn}
\begin{itemize}
\raggedright
  \item \emph{Say it:} \emph{pāip}
  \item \emph{Say it:} \emph{pāip}, once more
  \item \emph{Say it:} \emph{Rasoī kā pāip ṭūṭ gayā.} — The kitchen pipe has broken.
  \item \emph{Recall:} say the Hindi for to let drip, then the Hindi for to crack, to split, then say \emph{pāip} again
\end{itemize}

\begin{tcolorbox}[breakable,colback=teal!4,colframe=teal!35!black,title={Before you move on}]
What does \hi{पाइप} mean? (\textquotedblleft{}A pipe\textquotedblright{}.) Say it once more.
\end{tcolorbox}
\section[\texorpdfstring{\hi{छेद} (ched)}{छेद (ched)}]{\hi{छेद} (ched) — a hole}

\label{lesson:HI-C516-ched}



\begin{quote}
Before the new one: say the Hindi for to crack, to split, then the Hindi for a pipe.
\end{quote}

\subsection*{You'll want to know: \hi{छेद}}
\textbf{\hi{छेद}} — \emph{ched} — \textquotedblleft{}a hole\textquotedblright{}.

\begin{quote}
\hi{बाल्टी} \hi{में} \hi{एक} \hi{छोटा} \hi{छेद} \hi{है।}
\end{quote}

\begin{quote}
\emph{Bālṭī meṁ ek choṭā ched hai.}
\end{quote}

\begin{quote}
There is a small hole in the bucket.
\end{quote}

Masculine.

\begin{cousinweb}
\textbf{\hi{छेद}} (\emph{ched}) is a tatsama, taken whole from Sanskrit \emph{cheda}, \textquotedblleft{}cutting off; a section, an incision\textquotedblright{} (Monier-Williams); Hindi uses it for a hole.
\end{cousinweb}

\begin{grammarlens}[title={What you've built}]
One more word for describing what you lived and felt.
\end{grammarlens}

\subsection*{Your turn}
\begin{itemize}
\raggedright
  \item \emph{Say it:} \emph{ched}
  \item \emph{Say it:} \emph{ched}, once more
  \item \emph{Say it:} \emph{Bālṭī meṁ ek choṭā ched hai.} — There is a small hole in the bucket.
  \item \emph{Recall:} say the Hindi for to crack, to split, then the Hindi for a pipe, then say \emph{ched} again
\end{itemize}

\begin{tcolorbox}[breakable,colback=teal!4,colframe=teal!35!black,title={Before you move on}]
What does \hi{छेद} mean? (\textquotedblleft{}A hole\textquotedblright{}.) Say it once more.
\end{tcolorbox}
\section[\texorpdfstring{\hi{ढलकना} (ḍhalaknā)}{ढलकना (ḍhalaknā)}]{\hi{ढलकना} (ḍhalaknā) — to trickle down; to roll off}

\label{lesson:HI-C516-dhalakna}



\begin{quote}
Before the new one: say the Hindi for a pipe, then the Hindi for a hole.
\end{quote}

\subsection*{You'll want to know: \hi{ढलकना}}
\textbf{\hi{ढलकना}} — \emph{ḍhalaknā} — \textquotedblleft{}to trickle down; to roll off\textquotedblright{}.

\begin{quote}
\hi{छत} \hi{से} \hi{पानी} \hi{ढलककर} \hi{आँगन} \hi{में} \hi{गिरा।}
\end{quote}

\begin{quote}
\emph{Chat se pānī ḍhalakkar āṁgan meṁ girā.}
\end{quote}

\begin{quote}
Water ran off the roof and fell into the courtyard.
\end{quote}

Past (it …): \hi{वह} \hi{पानी} \hi{ढलक} \hi{गया}, \emph{vah pānī ḍhalak gayā}.

\begin{cousinweb}
\textbf{\hi{ढलकना}} (\emph{ḍhalaknā}) is a tadbhava, inherited from \emph{*ḍhalati}, \textquotedblleft{}bends over, falls\textquotedblright{}, a form CDIAL reconstructs from its descendants (CDIAL 5581). CDIAL gives the Hindi verb as \textquotedblleft{}to lean over, be spilt\textquotedblright{}.
\end{cousinweb}

\begin{grammarlens}[title={What you've built}]
One more word for describing what you lived and felt.
\end{grammarlens}

\subsection*{Your turn}
\begin{itemize}
\raggedright
  \item \emph{Say it:} \emph{ḍhalaknā}
  \item \emph{Say it:} \emph{ḍhalaknā}, once more
  \item \emph{Say it:} \emph{Chat se pānī ḍhalakkar āṁgan meṁ girā.} — Water ran off the roof and fell into the courtyard.
  \item \emph{Recall:} say the Hindi for a pipe, then the Hindi for a hole, then say \emph{ḍhalaknā} again
\end{itemize}

\begin{tcolorbox}[breakable,colback=teal!4,colframe=teal!35!black,title={Before you move on}]
What does \hi{ढलकना} mean? (\textquotedblleft{}To trickle down; to roll off\textquotedblright{}.) Say it once more.
\end{tcolorbox}
